\subsection{The normal basis theorem}

Let N be a Galois extension of K, with Galois group \(\Gamma\) . We identify \(\Gamma\) with the canonical basis of the group algebra \(\mathbf{K}^{(\Gamma)}\) (III, p.~446); then N may be considered as a left \(\mathbf{K}^{(\Gamma)}\) -module (III, p.~447, Example), so that

\[
u. x = \sum_ {\sigma \in \Gamma} a _ {\sigma} \sigma (x) \quad \text { for } \quad x \in \mathbf {N} \quad \text { and } \quad u = \sum_ {\sigma \in \Gamma} a _ {\sigma} \sigma \quad \text { in } \quad \mathbf {K} ^ {(\Gamma)}.
\]

If \(\mathbf{N}\) is of finite degree over \(\mathbf{K}\) , the group \(\Gamma\) is finite by Dedekind's theorem (V, p.~27, Cor. 2) and we can define the element \(t = \sum_{\sigma \in \Gamma} \sigma\) in \(\mathbf{K}^{(\Gamma)}\) ; then we have

\[
\operatorname{Tr} _ {N / K} (x) = \sum_ {\sigma \in \Gamma} \sigma (x),
\]

that is, \(\operatorname{Tr}_{\mathbf{N} / \mathbf{K}}(x) = t.x\) for all \(x \in \mathbf{N}\) .

Let us define an action on the right by \(\Gamma\) on \(\mathbf{N}\) by \(x^{\sigma} = \sigma^{-1}(x)\) . In a similar way we can consider the multiplicative group \(\mathbf{N}^{*}\) as a right \(\mathbf{Z}^{(\Gamma)}\) -module, the external law being written \((x, u) \mapsto x^u\) . For example, the notation \(x^{2\sigma + 3\tau + \pi}\) , where \(\sigma, \tau, \pi\) are elements of \(\Gamma\) , indicates the product \((x^\sigma)^2 \cdot (x^\tau)^3 \cdot x^\pi\) . If \(\mathbf{N}\) is of finite degree over \(\mathbf{K}\) and \(t = \sum_{\sigma \in \Gamma} \sigma\) as above, then we have \(\mathrm{N}_{\mathrm{N/K}}(x) = \prod_{\sigma \in \Gamma} x^\sigma\) , that is, \(\mathrm{N}_{\mathrm{N/K}}(x) = x^t\) for all \(x \in \mathbf{N}^*\) .

Suppose henceforth that N is of finite degree over K. Given \(x \in N\) , for \(\{x\}\) to be a basis of the \(\mathbf{K}^{(\Gamma)}\) -module N it is necessary and sufficient that the family \((\sigma(x))_{\sigma \in \Gamma}\) should be a basis of N over K. Such a basis is called a normal basis of N over K.

THEOREM 6. --- Let N be a Galois extension of finite degree over K and let \(\Gamma\) be the Galois group of N over K. Then there exists a normal basis of N over K; in other words, the \(\mathbf{K}^{(\Gamma)}\) -module N is free of rank 1.

We shall give two proofs of this statement. The first uses the following lemma which will be proved in Chapter VIII (§ 2, No.~5).

* Lemma 3. --- Let A be a K-algebra, \(\mathbf{M}_1\) and \(\mathbf{M}_2\) two A-modules of finite rank over K and suppose that there exists an extension L of K such that the modules \(\mathrm{L} \otimes_{\mathrm{K}} \mathrm{M}_1\) and \(\mathrm{L} \otimes_{\mathrm{K}} \mathrm{M}_2\) over the ring \(\mathrm{L} \otimes_{\mathrm{K}} \mathrm{A}\) are isomorphic. Then the A-modules \(\mathbf{M}_1\) and \(\mathbf{M}_2\) are isomorphic.

We shall apply Lemma 3 in the case where \(\mathbf{A} = \mathbf{K}^{(\Gamma)}\) , \(\mathbf{M}_1 = \mathbf{N}\) , \(\mathbf{M}_2 = \mathbf{A}_s\) and \(\mathbf{L} = \mathbf{N}\) . By the Cor. of V, p.~64 there exists a K-isomorphism \(\varphi\) of \(\mathbf{N} \otimes_{\mathbf{K}} \mathbf{N}\) onto \(\mathbf{N} \otimes_{\mathbf{K}} \mathbf{K}^{(\Gamma)}\) which maps \(x \otimes y\) to \(\sum x \sigma^{-1}(y) \otimes \sigma\) . It is clear that \(\varphi\) is an isomorphism of \(\mathbf{N} \otimes_{\mathbf{K}} \mathbf{K}^{(\Gamma)}\) -modules and so the theorem follows from Lemma 3. *

For the second proof we shall use the following proposition :

PROPOSITION 10. --- Let \(x \in \mathbb{N}\) , then for \(\{x\}\) to form a basis of the \(K^{(\Gamma)}\) -module \(\mathbb{N}\) it is necessary and sufficient that \(\det(\sigma\tau(x))_{\sigma,\tau \in \Gamma} \neq 0\) .

Since \(\mathbf{K}^{(\Gamma)}\) and \(\mathbf{N}\) have the same dimension over \(\mathbf{K}\) , to say that \(\{x\}\) is a basis of \(\mathbf{N}\) over \(\mathbf{K}^{(\Gamma)}\) means that the mapping \(a \mapsto ax\) of \(\mathbf{K}^{(\Gamma)}\) into \(\mathbf{N}\) is injective. This in turn means that the mapping \(b \mapsto b(1 \otimes x)\) of \(\mathbf{N} \otimes_{\mathbf{K}} \mathbf{K}^{(\Gamma)}\) into \(\mathbf{N} \otimes_{\mathbf{K}} \mathbf{N}\) is injective (II, p.~306, Prop. 14). Now there is an \(\mathbf{N} \otimes_{\mathbf{K}} \mathbf{K}^{(\Gamma)}\) -module isomorphism of \(\mathbf{N} \otimes_{\mathbf{K}} \mathbf{N}\) onto \(\mathbf{N} \otimes_{\mathbf{K}} \mathbf{K}^{(\Gamma)}\) which maps \(1 \otimes x\) to \(\sum_{\sigma} \sigma^{-1}(x) \otimes \sigma\) . It follows that \(\{x\}\) is a basis

of N over \(\mathbf{K}^{(\Gamma)}\) if and only if, for every non-zero family of elements \((n_{\tau})_{\tau\in\mathrm{T}}\) of N we have \(\left(\sum n_{\tau}\otimes\tau\right)\left(\sum\sigma^{-1}(x)\otimes\sigma\right)\neq0\) . But this last relation means that there exists \(\sigma\in\Gamma\) such that \(\sum_{\tau}n_{\tau}\sigma^{-1}\tau(x)\neq0\) , whence the proposition.

\begin{enumerate}
\def\labelenumi{\Alph{enumi})}
\tightlist
\item
  Suppose that K is infinite; the mapping \(x \mapsto \det(\sigma\tau(x))\) of N into N is a polynomial mapping over K (IV, p.~54). By extension of scalars from K to N we get a corresponding mapping for the vector N-space \(\mathbf{N} \otimes_{\mathbf{K}} \mathbf{N}\) and we have just seen that the latter is not identically zero (because \(\mathbf{N} \otimes_{\mathbf{K}} \mathbf{N}\) is free of rank 1 over \(\mathbf{N} \otimes_{\mathbf{K}} \mathbf{K}^{(\Gamma)}\) ). So there exists \(x \in \mathbf{N}\) such that \(\det(\sigma\tau(x)) \neq 0\) (IV, p.~18, Th. 2); more generally, from the same reference we have:
\end{enumerate}

PROPOSITION 11. --- Suppose that K is infinite and let P:N→K be a polynomial mapping which is non-zero on K. There exists x∈N such that P(x)≠0 and \{x\} is a basis of N over K \(^{(\Gamma)}\) .

\begin{enumerate}
\def\labelenumi{\Alph{enumi})}
\setcounter{enumi}{1}
\tightlist
\item
  Suppose that K is finite. By Prop. 4 (V, p.~95) \(^{1}\) every extension of finite degree over K has a cyclic Galois group. We shall therefore more generally consider the case where the group \(\Gamma\) is cyclic of order n; we denote by \(\gamma\) a generator of \(\Gamma\) .
\end{enumerate}

The following lemma is a particular case of more general results proved in Chapter VII. The ring A is either the ring Z of rational integers or the ring K{[}X{]} of polynomials over the field K.

Lemma 4. --- Let M be a torsion A-module generated by a finite number \(x_{1}, \ldots, x_{h}\) of elements; then there exists an element x of M whose annihilator (II, p.~219) is equal to the annihilator of M.

In both cases A is an integral domain and every ideal of A is principal. When A = Z (resp. A = K{[}X{]}), we denote by P the set of prime numbers (resp. the set of irreducible monic polynomials in K{[}X{]}). For every element \(a \neq 0\) of A there exists then an invertible element u of A and a family \((v_{p}(a))_{p \in \mathcal{P}}\) , with finite support, of positive integers such that \(a = u \prod_{p \in \mathcal{P}} p^{v_{p}(a)}\) and u and the integers

\(v_{p}(a)\) are uniquely determined (I, p.~51 and IV, p.~13, Prop. 13).

Let \(\mathfrak{a}_i\) be the annihilator of \(x_i\) (for \(1 \leqslant i \leqslant h\) ) and \(\mathfrak{a}\) the annihilator of \(\mathbf{M}\) ; let \(a_1, \ldots, a_h\) , \(a\) be non-zero elements of \(\mathbf{A}\) such that \(\mathfrak{a}_i = \mathbf{A}a_i\) and \(\mathfrak{a} = \mathbf{A}a\) ; since \(\mathfrak{a} = \mathfrak{a}_1 \cap \ldots \cap \mathfrak{a}_h\) , it follows from what has been said that

\[
v _ {p} (a) = \sup _ {1 \leqslant i \leqslant h} v _ {p} (a _ {i}) \quad \text { for   all } p \in \mathscr {P}.\tag{10}
\]

Let us write \(a\) in the form \(up_1^{n(1)} \ldots p_r^{n(r)}\) , with \(p_1, \ldots, p_r\) distinct in \(\mathcal{P}\) , \(n(1) > 0, \ldots, n(r) > 0\) and \(u\) an invertible element of \(A\) . Let \(j = 1, \ldots, r\) ; by (10) there exists an integer \(c(j)\) such that \(1 \leqslant c(j) \leqslant h\) and \(v_{p_j}(a_{c(j)}) = n(j)\) ; there exists \(b_j\) in \(A\) with \(a_{c(j)} = p_j^{n(j)}b_j\) and the element \(y_j = b_jx_{c(j)}\) has as annihilator the ideal \(Ap_j^{n(j)}\) .

Let us show that the annihilator b of \(y = y_{1} + \cdots + y_{r}\) is equal to the annihilator a of M. In any case we have \(a \subset b\) , so b is of the form \(Ap_{1}^{m(1)} \ldots p_{r}^{m(r)}\) with \(0 \leqslant m(j) \leqslant n(j)\) for \(1 \leqslant j \leqslant r\) . If we had \(a \neq b\) , there would exist an integer j such that \(1 \leqslant j \leqslant r\) and \(m(j) < n(j)\) , and hence \(d_{j} = a/p_{j}\) would annihilate y. Now we have \(d_{j}y_{k} = 0\) for \(k \neq j\) , whence we obtain \(d_{j}y_{j} = 0\) ; but the annihilator of \(y_{j}\) is \(Ap_{j}^{m(j)}\) and \(d_{j}\) is not a multiple of \(p_{j}^{m(j)}\) . So the hypothesis \(a \neq b\) is absurd.

We shall apply Lemma 4 to the case where \(\mathbf{A}\) is the polynomial ring \(\mathbf{K}[\mathbf{X}]\) and \(\mathbf{M}\) the abelian group \(\mathbf{N}\) with the external law defined by \(a \cdot x = \sum_{k=0}^{\infty} c_k \gamma^k(x)\) for \(a = \sum_{k=0}^{\infty} c_k \mathbf{X}^k\) in \(\mathbf{K}[\mathbf{X}]\) and \(x \in \mathbb{N}\) . Let \(\mathfrak{a}\) be the annihilator of \(\mathbf{M}\) , then \(\gamma^n = 1\) , hence the polynomial \(\mathbf{X}^n - 1\) lies in \(\mathfrak{a}\) . Let \(\mathbf{F} \in \mathfrak{a}\) ; by IV, p.~11, Cor. there exist elements \(c_0, c_1, \ldots, c_{n-1}\) of \(\mathbf{K}\) and \(\mathbf{G} \in \mathbf{K}[\mathbf{X}]\) such that

\[
\mathrm{F} (\mathbf {X}) = c _ {0} + c _ {1} \mathbf {X} + \dots + c _ {n - 1} \mathbf {X} ^ {n - 1} + (\mathbf {X} ^ {n} - 1) \mathrm{G} (\mathbf {X}).\tag{11}
\]

Thus we have \(c_{0} + c_{1}\gamma + \cdots + c_{n-1}\gamma^{n-1} = 0\) in \(\operatorname{Hom}_{\mathbf{K}}(\mathbf{N}, \mathbf{N})\) and since the automorphisms 1, \(\gamma\) , \(\gamma^{2}\) , \ldots, \(\gamma^{n-1}\) of N are distinct, Dedekind's theorem (V, p.~27, Cor. 2) implies that \(c_{0} = c_{1} = \cdots = c_{n-1} = 0\) . Finally, we have

\[
\mathrm{F} (\mathbf {X}) = \left(\mathbf {X} ^ {n} - 1\right) \mathrm{G} (\mathbf {X}), \quad \text { that   is }, \quad \mathfrak {a} = \left(\mathbf {X} ^ {n} - 1\right) \mathrm{K} [ \mathbf {X} ].
\]

By Lemma 4 there exists an element \(x\) of \(\mathbf{N}\) whose annihilator in \(\mathbf{K}[\mathbf{X}]\) is equal to \((\mathbf{X}^n - 1)\mathbf{K}[\mathbf{X}]\) . Since the monomials 1, \(\mathbf{X}, \ldots, \mathbf{X}^{n-1}\) form a basis of a vector subspace of \(\mathbf{K}[\mathbf{X}]\) supplementary to \((\mathbf{X}^n - 1)\mathbf{K}[\mathbf{X}]\) (IV, p.~11, Cor.), the elements \(x, \gamma(x), \ldots, \gamma^{n-1}(x)\) of \(\mathbf{N}\) are linearly independent over \(\mathbf{K}\) . Since \([\mathbf{N}:\mathbf{K}] = n(\mathbf{V},\mathbf{p}.66,\mathbf{Th}.3)\) , the sequence \((x,\gamma(x),\ldots,\gamma^{n-1}(x))\) is therefore a (normal) basis of \(\mathbf{N}\) over \(\mathbf{K}\) .

